\subsection*{Widening the scope: cooperative group and mma}

A wider cooperating group uses \co@: group-4@ while retaining the same fill, load, multiply, and store operations. The example also changes tile dimensions and staging buffers to suit that scope. The shared syntax makes the two versions comparable, but a change in cooperation width still requires compatible tile geometry and target support.

\begin{figure}[htbp]\centering
\includegraphics[width=\linewidth,height=0.56\textheight,keepaspectratio]{f06-fig4_sm86_vs_sm90_light.png}
\caption{The : group and : group-4 scopes use different cooperation widths while preserving the fill, load, multiply, and store interface.}
\end{figure}

\begin{lstlisting}[language=croqtile,caption={Widening the scope: cooperative group and mma}]
#define TILE_M 128
#define TILE_N 128
#define TILE_K 128
#define WARP_M 64
#define WARP_N 64
#define WARP_K 16
__co__ void matmul_wgmma(global f16 [M, K] lhs, global f16 [N, K] rhs, global f16 [M, N] output) {
  parallel {block_m, block_n} by [cdiv(M, WARP_M), cdiv(N, WARP_N)] : block {
    shared f16 [WARP_M, TILE_K] lhs_load_s;
    shared f16 [WARP_N, TILE_K] rhs_load_s;
    mc = mma.fill.f16 0.0f;
    foreach {iv_k} in [cdiv(K, TILE_K)] {
      dma.copy lhs.subspan(WARP_M, TILE_K).at(block_m, iv_k) => lhs_load_s;
      dma.copy rhs.chunkat(block_n, iv_k) => rhs_load_s;
      foreach {iv_warp} in [cdiv(TILE_K, WARP_K)] {
        parallel p by 1 : group-4 {
          ma = mma.load lhs_load_s.chunkat(_, iv_warp);
          mb = mma.load rhs_load_s.chunkat(_, iv_warp);
          mma.row.row mc, ma, mb;
        }
      }
    }
    shared f16 [WARP_M, WARP_N] output_s;
    mma.store mc, output_s;
    dma.copy output_s => output.subspan(WARP_M, WARP_N).at(block_m, block_n);
  }
}
\end{lstlisting}

The sequence remains fill, load, multiply, and store; the space specifier determines the cooperation scope.

The fill suffix chooses accumulator precision. This listing uses \co@mma.fill.f16 0.0f@, so the initial single-precision literal is converted for an FP16 accumulator. An FP32 accumulator is a separate choice, expressed with \co@.f32@, and provides a wider range and greater precision for partial sums. The operand type alone should not be used to infer what an explicitly annotated accumulator stores.

The expression \co@chunkat(_, iv_warp)@ leaves the row dimension whole and subdivides only K. The underscore means that no partition is applied to that dimension.

The underscore in the operand selection leaves one dimension intact while the loop divides K into MMA slices. Each slice contributes to the same accumulator. This lets the outer staging tile be larger than the fragment consumed by one matrix operation without changing the output tile being accumulated.
